\chapter{Introduction}

\section{Contexte du projet}

Le projet DEUS Dashboard est une application web permettant de visualiser des donnees
astronomiques fournies par les APIs de la NASA. L'application affiche l'image astronomique
du jour (APOD), les evenements naturels terrestres (EONET), ainsi qu'une carte interactive
de la Lune.

L'objectif de ce projet est de mettre en \oe{}uvre une methodologie DevOps complete
pour deployer cette application sur Amazon Web Services (AWS), en demontrant la maitrise
des outils et pratiques essentiels du domaine.

\section{Objectifs}

\begin{itemize}
    \item Conteneuriser l'application avec Docker
    \item Provisionner l'infrastructure AWS avec Terraform (Infrastructure as Code)
    \item Configurer les serveurs avec Ansible (Configuration Management)
    \item Automatiser les deploiements avec GitHub Actions (CI/CD)
    \item Mettre en place un monitoring avec CloudWatch
    \item Securiser l'infrastructure (IAM, VPC, SSM Parameter Store)
\end{itemize}

\section{Technologies utilisees}

\begin{table}[H]
    \centering
    \begin{tabularx}{\textwidth}{lX}
        \toprule
        \textbf{Categorie} & \textbf{Technologies} \\
        \midrule
        Frontend & Angular 20, TypeScript, Leaflet \\
        Backend & Node.js, Express 5, TypeScript \\
        Base de donnees & PostgreSQL 16 (RDS), DynamoDB \\
        Conteneurisation & Docker, Docker Compose \\
        Infrastructure & Terraform, AWS (VPC, ECS, ALB, S3, CloudFront) \\
        Configuration & Ansible \\
        CI/CD & GitHub Actions \\
        Monitoring & CloudWatch (Logs, Metriques, Alarmes, Dashboard) \\
        Securite & IAM, Security Groups, SSM Parameter Store \\
        \bottomrule
    \end{tabularx}
    \caption{Stack technologique du projet}
\end{table}
